\chapter{Introdução}

\section{Contextualização}

O uso de baterias de íons de lítio em sistemas embarcados, mobilidade elétrica, armazenamento
de energia e equipamentos portáteis tem crescido de forma significativa nas últimas décadas.
Com esse crescimento, aumenta também a necessidade de sistemas de gerenciamento de baterias
(\textit{Battery Management Systems} -- BMS) capazes de monitorar, proteger e operar os packs
de forma segura \cite{andrea2010,plett2015v1}.

\section{Problema de pesquisa}

Em packs multicélula, como no caso de uma bateria \textit{4S}, a simples leitura da tensão total
do pack não é suficiente para garantir segurança operacional. É necessário monitorar cada célula
individualmente, acompanhar temperatura, corrente e estados de falha, além de aplicar estratégias
de proteção e balanceamento.

\section{Objetivo geral}

Desenvolver a proposta de uma BMS para bateria de lítio em configuração \textit{4S}, utilizando
um ESP32-S3 como unidade de supervisão e um \textit{Analog Front End} dedicado para medição e
proteção primária do pack.

\section{Objetivos específicos}

\begin{itemize}
  \item identificar os requisitos de proteção de uma BMS para pack 4S;
  \item analisar o papel do AFE em sistemas de gerenciamento de baterias;
  \item propor a arquitetura de hardware da solução;
  \item estruturar a máquina de estados e a lógica de proteção do firmware;
  \item organizar uma base de documentação e configuração para futura implementação física.
\end{itemize}

\section{Justificativa}

A relevância do tema se apoia na necessidade de segurança, robustez e rastreabilidade em sistemas
com células de lítio. Uma BMS mal projetada pode causar degradação precoce do pack, perda de
desempenho e, em casos severos, risco elétrico e térmico.

\section{Metodologia}

Este trabalho adota uma abordagem de pesquisa aplicada, de natureza tecnológica, combinando:

\begin{itemize}
  \item revisão bibliográfica sobre BMS, AFE e proteção de packs de lítio;
  \item análise de arquiteturas e componentes de mercado;
  \item modelagem da solução para uma BMS 4S;
  \item estruturação inicial do firmware e da documentação técnica.
\end{itemize}

\section{Estrutura do trabalho}

Este trabalho está organizado em sete capítulos. Após esta introdução, o Capítulo 2 apresenta
a fundamentação teórica. O Capítulo 3 descreve os requisitos e a metodologia. O Capítulo 4
apresenta a arquitetura proposta. O Capítulo 5 detalha o desenvolvimento da solução. O
Capítulo 6 discute os resultados e limitações. Por fim, o Capítulo 7 apresenta a conclusão.
